\section{Methodology}\label{sec:method}

\subsection{The overnight-stay networks}

The Ministry of Transport publishes, for each night, an estimate of how many residents of each zone slept in each other zone \citep{mitms2024methodology}. According to the ministry's methodology, the estimate comes from the records of a mobile-phone operator: each phone is given a home zone, where it usually spends the night, and a zone for each night; the counts are then scaled up to the resident population. I use the 2,435 zones of mainland Spain, so that the network covers one connected territory, and the six months of 2024 in the data used here: February, March, August, October, November and December. A file dated $d$ holds the night that ends on the morning of $d$; I always name a night by both dates, as in ``14--15 August''.

For each night I build a directed network. A link $i\to j$ exists when at least 10 residents of $i$ sleep in $j$, and carries the estimated number of people as its weight. Stays inside one's own zone are not links. Below 10 people the estimates are noisy and the many small links would make every network almost complete; with the threshold, the networks keep 69--83\,\% of the people who sleep away from home and 0.9--3.4\,\% of all possible links, well below the 5\,\% limit of sparse networks. Because the threshold keeps more links on busy nights, every structural comparison is repeated on a second version in which each night keeps the same number of links as the ordinary night, 35,812, chosen by weight (Table~\ref{tab:basic_full}).

Three nights are the main networks: an ordinary weekday (13--14 February), the summer holiday (14--15 August) and All Saints (31 October--1 November). Any one night could be unusual, so each result is repeated on its \emph{sister nights}: the 13 Tuesday, Wednesday and Thursday nights of February, the 9 nights ending on 10--18 August and the 3 nights ending on 1--3 November. A holiday result is only reported if it has the same sign in every pair of one holiday sister and one February sister (117 pairs for summer, 39 for All Saints), and the main nights are repeated in 2022 and, except All Saints, in 2023.

Table~\ref{tab:basic} and Figure~\ref{fig:basic} give the basic properties of the three networks, computed on their undirected version (zones $i$ and $j$ are linked if either sends people to the other). All three are small worlds: any zone reaches any other in about two steps, close to the value $\ell^*=\ln N/\ln\langle k\rangle$ expected in a random graph \citep{zhou2007structural}, and their clustering is 0.55--0.59, against 0.01--0.03 in random graphs with the same number of nodes and links. Their degree distributions are heavy-tailed but not power laws: a likelihood-ratio test \citep{clauset2009power} cannot tell a power law from a log-normal on the ordinary and All Saints nights and prefers the log-normal in summer. All three are disassortative: well-connected zones link mostly to poorly connected ones.

\input{../tables/table1}

\begin{figure}[t]
\centering
\includegraphics[width=\linewidth]{fig_basic_properties.pdf}
\caption{\textbf{Basic properties of the three networks.} \textbf{a}, Share of zones with degree at least $k$. \textbf{b}, Mean degree of a zone's neighbours, $k_{nn}$, against its degree, in logarithmic bins. \textbf{c}, Rich-club connectivity $\phi$ against the normalised rank $r/N$ of the best-connected zones. Undirected networks with links of at least 10 people; dotted lines show degree-preserving rewired networks (mean of 10 in b and 50 in c).}
\label{fig:basic}
\end{figure}

\subsection{The roots network}

Where residents were born comes from the 10\,\% sample of the 2024 census \citep{ine2025censo}, which gives, for each person, the municipality of residence (if it has more than 10,000 inhabitants), the province of birth, the size of the municipality of birth, and the year of arrival in the current province. I keep the 675 municipalities over 10,000 inhabitants that form a zone of their own and lie in one of the 44 provinces with at least two of them; this is needed to compare neighbouring municipalities, below. For each municipality $c$ and each other mainland province $p$, the roots network counts the residents of $c$ born in $p$, in total and split into groups: by arrival (before 1981, which covers the exodus generation, and 2016--2023) and by birthplace (a village of 2,000 inhabitants or fewer, or a town of more than 10,000). I call the 675 municipalities \emph{cities}, although many are small towns. For the total pull of the roots, without groups, I also use the full census table of residents by province of birth, which covers 137 large mainland cities; 112 of them lie in the 22 provinces with at least two, which identify the comparison.

\subsection{Methods}

\paragraph{Structure against degree-preserving null models (T1, T2).} On holiday the network is much denser, and most of its properties could change for that reason alone. To separate what the new degrees force from what they do not, I compare each network with 50 random networks in which every zone keeps its degree (each made with 10 random link swaps per link), and report the real value next to their mean. For degree mixing I follow \citet{zhou2007structural}: rewiring the network towards maximal assortativity and maximal disassortativity, while keeping every degree and a single connected component, gives the range of the assortativity $r$ that the degrees allow, and shows where the real network sits in it. The rich club is measured as $\phi(r)$ against the rank $r$ of the best-connected zones \citep{zhou2004richclub}, first in raw form and then against the rewired networks, with the warning of \citet{zhou2007structural} that a ratio close to one does not mean there is no rich club when the null networks are themselves dominated by hubs. Second-order mixing $R_{\mathrm{avg}}$ \citep{zhou2017second}, which asks whether the neighbours of linked zones are alike, is compared with rewired networks that keep both the degrees and the rich club. The appendix gives the bitmap of the sorted adjacency matrix (BOSAM) of each night \citep{guo2009fingerprint}.

\paragraph{The pull of the roots (T3).} The test asks whether a city sends more of its residents to the provinces where more of them were born. The difficulty is that a province can attract people for many reasons: it may be close, cheap, or have a coast. I remove these by comparing two cities of the same province that send people to the same destination province: everything that belongs to the pair of provinces cancels, and the comparison is between cities with more or fewer residents born there. Formally, for city $c$ in province $o(c)$, destination province $p\neq o(c)$ and night $t$, the number of residents of $c$ who sleep in $p$, $y_{cpt}$, is modelled by Poisson pseudo-maximum likelihood,
\begin{equation}
\mathrm{E}[y_{cpt}] = \exp\big(\beta_t \log s_{cp} + \alpha_{ct} + \gamma_{o(c)pt} + f_t(d_{cp}) + \lambda_{pt}\,\log I_c\big),
\label{eq:ppml}
\end{equation}
where $s_{cp}$ is the share of the city's residents born in $p$ (a half is added to every count, so that empty cells can be used), $\alpha$ are city effects, $\gamma$ are origin-province $\times$ destination-province effects, $f_t$ is a quadratic in the log distance from the city to the centre of $p$, and the last term lets the income per person of the city, $I_c$, affect each destination differently. The coefficient $\beta_t$ is the \emph{pull of the roots}: a city with twice the share of residents born in $p$ sends $2^{\beta_t}$ times as many people there. Two nights are fitted together, so that the change $\delta=\beta_{\text{holiday}}-\beta_{\text{ordinary}}$ is estimated and tested directly. To find out which residents carry the pull, $\log s_{cp}$ is replaced by the shares of each group, entered together: arrived before 1981 and arrived in 2016--2023, or born in a village and born in a town. The model is fitted separately for stays in inland and in coastal zones of $p$. With only 44 provinces, standard errors are unreliable, so every $p$-value and interval comes from a wild cluster bootstrap by province of origin \citep{kline2012score}, with 9,999 draws.

\paragraph{The emptied villages (T4).} For each zone I measure how much population it lost: one minus the ratio of its 2021 population to the largest population it had between 1900 and 1960. Zones are split into five groups of about 475 by this loss. A zone is a \emph{net receiver} on a night when more people from elsewhere sleep in it than residents of the zone sleep elsewhere, using all flows, not only links above the threshold. I compare the share of net receivers in each group across nights and test the change by regression, holding fixed zone population, the number of municipalities in the zone, area, the coast, the distance to the nearest city of 100,000 people and to the coast, and the province. Three zones with gaps in the data and 56 without population history are excluded, leaving 2,376.

\paragraph{Evaluation.} Besides the sister nights and the other years, the roots results are checked by leaving out one province at a time and against an October baseline instead of February, and the village results on zones made of a single municipality, to rule out an effect of grouping small municipalities into one zone. Section~\ref{sec:results} reports which results pass these checks and which do not.
